\documentclass[11pt]{article}
\usepackage[margin=1in]{geometry}
\usepackage{amsmath,amssymb,amsthm,mathtools}
\usepackage[T1]{fontenc}
\usepackage{lmodern,microtype,enumitem,xurl}
\usepackage[colorlinks=true,linkcolor=blue,citecolor=blue,urlcolor=blue]{hyperref}
\hypersetup{pdftitle={Arithmetic formulas for low-product tanh networks and a Bayesian learning converse},pdfauthor={ChatGPT-6 Astra; Kaarel Haenni; Dmitry Vaintrob}}
\newtheorem{theorem}{Theorem}[section]
\newtheorem{lemma}[theorem]{Lemma}
\newtheorem{proposition}[theorem]{Proposition}
\newtheorem{corollary}[theorem]{Corollary}
\newtheorem{definition}[theorem]{Definition}
\theoremstyle{remark}\newtheorem{remark}[theorem]{Remark}
\newcommand{\R}{\mathbb R}
\newcommand{\C}{\mathbb C}
\newcommand{\E}{\mathbb E}
\newcommand{\Prb}{\mathbb P}
\newcommand{\norm}[1]{\left\lVert#1\right\rVert}
\newcommand{\poly}{\operatorname{poly}}
\newcommand{\PNC}{\mathsf{PNC}}
\newcommand{\NC}{\mathsf{NC}}
\newcommand{\AC}{\mathsf{AC}}
\newcommand{\Lpoly}{\mathsf L/\mathrm{poly}}
\newcommand{\Ppoly}{\mathsf P/\mathrm{poly}}
\newcommand{\PF}{P_F}
\newcommand{\MAJ}{\operatorname{MAJ}}
\newcommand{\sign}{\operatorname{sign}}
\newcommand{\Rea}{\operatorname{Re}}
\newcommand{\Ima}{\operatorname{Im}}
\title{Arithmetic formulas for low-product tanh networks\
and a Bayesian learning converse}
\author{ChatGPT-6 Astra \and Kaarel H\"anni \and Dmitry Vaintrob\thanks{Authors are listed in alphabetical order.}}
\date{Research draft --- 25 October 2026}
\begin{document}
\maketitle
\begin{abstract}
We prove that every Boolean function with a globally unit-margin tanh
representation of Frobenius product $P$ has a logarithmic-depth integer
arithmetic circuit of size polynomial in $d+P$ whose output sign equals
the function. The original network may have arbitrary finite depth,
widths, and real coefficients. In particular, polynomial product implies
nonuniform $\PNC^1$, hence logarithmic space with polynomial advice.
The analytic ingredient is an exact Fourier--Poisson estimator for tanh
that multiplies expected arithmetic work times second moment by only a
universal constant. Cost-aware neuron sampling makes this potential
propagate according to squared Frobenius norms. A logarithmic-depth tanh
normal form, finite-set derandomization, coefficient rounding, and
real-affine formula balancing complete the proof.

For the original Gaussian Bayesian network conditioned on a global
absolute output margin and exact training labels, an established
posterior representation ceiling then gives a converse: a
distribution-free guarantee after $n$ samples, with confidence-failure
mass below any fixed constant less than one half, forces a small
arithmetic-sign representation of the target. This applies to the
strong-prior limit and every fixed Gaussian exponent $\gamma>1$ at the
specified sufficiently large widths. It handles the full posterior
mixture. The endpoint $\gamma=1$ is not covered by this strengthening.
No unconditional exponential-depth sample lower bound or general
polynomial-depth learning upper bound follows. The new arguments have
undergone mathematical review but have not been formalized in Lean.
\end{abstract}

\section{Introduction and results}

How much sequential computation can a neural network hide behind a
small product of layer norms? This question matters for strong Gaussian
Bayesian priors: after conditioning on data, most posterior functions
can have low-product representatives even when the sampled parameters
themselves are numerous and not small. A structural theorem about these
functions can therefore turn a learning guarantee into a necessary
computational condition on the target.

For tanh networks, we show that polynomial Frobenius product restricts
the represented Boolean function to a class between Boolean logarithmic
depth and logarithmic space. This conclusion is stronger than merely
exhibiting a small Boolean circuit. It also differs from simulating the
given network layer by layer: the original network can have unbounded
finite depth and arbitrary real weights.

The proof has two parts. A geometric covering argument reduces the
function to an ordinary-tanh network of logarithmic depth and polynomial
resources. We include this normal form and its required width
compression in the appendices. We then replace the remaining nonlinear
computation by a small collection of randomized arithmetic expressions.
An exact Fourier representation and Poisson random products prevent
second-moment losses from escalating with the moment order. The useful
quantity is the product of expected work and second moment. It grows by
a constant times the squared layer Frobenius norm, so logarithmic
retained depth leaves polynomial cost. The probabilistic method selects
expressions correct on the whole Boolean cube. A careful finite-bit
conversion makes their coefficients legitimate circuit constants.

\subsection{Conventions}

Write $X_d=\{-1,1\}^d$ and augment the input by one constant coordinate.
A depth-$D$ network has $D$ weighted layers,
\[
 h^0(x)=(1,x),\qquad h^s(x)=\tanh(W^s h^{s-1}(x))\ (s<D),
 \qquad z(x)=W^D h^{D-1}(x).
\]
The readout is scalar and linear. Hidden widths are arbitrary finite
integers. There are no later free biases; a constant must be carried by
ordinary hidden neurons. Put
\[
 \PF(W)=\prod_{s=1}^D\norm{W^s}_F.
\]
A representation of $f:X_d\to\{-1,1\}$ has \emph{global unit margin}
if $f(x)z(x)\ge1$ for every $x\in X_d$. Purely affine representations
($D=1$) are allowed.

An integer arithmetic circuit uses binary $+$ and multiplication gates,
input bits, and constants $-1,0,1$. The nonuniform class $\PNC^1$
consists of signs of polynomial-size, logarithmic-depth such circuits:
the accepting output is strictly positive. This is an arithmetic sign
class, not Boolean $\NC^1$. Standard results give
\[
 \NC^1\subseteq\PNC^1\subseteq\Lpoly
 \quad\hbox{in the nonuniform setting},
 \qquad \AC^0(\PNC^1)=\PNC^1.
\]
Here and below the notation $\PNC^1$ is explicitly nonuniform unless a
uniform evaluator is mentioned. The last display uses the corresponding
nonuniform circuit form of the standard closure theorem
\cite{datta}; the advice consists of the selected finite circuits.

\begin{theorem}[Low product gives an arithmetic-sign circuit]
\label{thm:main}
There are universal constants $a,c>0$ such that any $f:X_d\to\{-1,1\}$
with a global unit-margin tanh representation of product at most $P$
has an integer arithmetic circuit of size at most $(d+P+2)^a$ and
depth at most $c\log(d+P+2)$ whose output is positive exactly when
$f(x)=1$. Its only constants are $-1,0,1$.
Consequently, any family with $P=\poly(d)$ belongs to nonuniform
$\PNC^1\subseteq\Lpoly$.
\end{theorem}

The theorem is existential. It does not say that the original real
parameters are computable, that the circuit can be found efficiently,
or that the original network can be evaluated in logarithmic space
from an arbitrary real encoding. It gives another finite representation
of the same Boolean function.

\begin{theorem}[Bayesian consequence, informal]
\label{thm:intro-converse}
For the Gaussian posterior specified in Section~\ref{sec:posterior},
fix either an exponent $\gamma>1$ and a sufficiently large finite
learner width, or a sufficiently large fixed width followed by the
strong-prior limit. Suppose that for every input distribution, $n$
labeled samples have positive probability of giving posterior
probability greater than $2/3$ to the correct label outside a set of
mass at most $\varepsilon<1/2$. Then the target has an arithmetic-sign
circuit of size polynomial in $d+n+(1-2\varepsilon)^{-1}$ and
logarithmic depth in the same quantities. The polynomial is independent
of learner depth and width. Precise width and limit quantifiers appear
in Theorems~\ref{thm:ceiling} and~\ref{thm:converse}.
\end{theorem}

If a single distribution-free sample bound polynomial jointly in the
resources of arbitrary bounded-weight deep tanh teachers and learner
depth held in either of these regimes, it would follow that
$\Ppoly$ is contained in nonuniform $\PNC^1$, and hence
$\Ppoly=\Lpoly$; see Corollary~\ref{cor:collapse}. This is a
conditional obstruction, not a proved separation. The question whether
unrestricted tanh Bayesian learning has polynomial or necessarily
exponential depth dependence remains open.

\subsection{Relation to arithmetic neural computation}

Merrill, Jiang, Li, Lin, and Sabharwal relate linear recurrent networks
to logarithmic-depth arithmetic circuits and distinguish them from
nonlinear recurrent architectures \cite{merrill}. Our setting is a
different restriction: the network is genuinely nonlinear and
feedforward, and neither its original depth nor its width is bounded.
The hypothesis is low Frobenius product together with a global Boolean
margin. The Fourier--Poisson argument constructs a new arithmetic
representation from this norm constraint; it does not replace the
nonlinear recurrence by an exact linear recurrence. The cited recurrent
results therefore do not supply the theorem here. The arithmetic-class
closure and logspace evaluation facts themselves are established
results of counting complexity \cite{datta,caussinus}.

\section{An exact random arithmetic expression for tanh}
\label{sec:estimator}

Expressions in this section may use complex arithmetic internally and
real or imaginary part at a unary gate. Their final outputs are real.
Cost counts leaves and arithmetic operations. Every child call uses
independent fresh randomness, and the random expression's shape is
independent of its input. Sampling exact real constants is a
nonuniform representation device; no computational cost bound for an
exact sampling algorithm is asserted.

\begin{lemma}[Fourier identity]
\label{lem:fourier}
For real $u$,
\[
 \tanh u=\int_0^\infty k(\omega)\sin(\omega u)\,d\omega,
 \qquad k(\omega)=\frac1{\sinh(\pi\omega/2)},
 \qquad \int_0^\infty\omega k(\omega)\,d\omega=1.
\]
The first integral is absolutely convergent.
\end{lemma}
\begin{proof}
The expansion $k(\omega)=2\sum_{j\ge0}e^{-(j+1/2)\pi\omega}$,
combined with integration of sine against an exponential, gives
\[
 \int_0^\infty k(\omega)\sin(\omega u)\,d\omega
 =2\sum_{j\ge0}\frac{u}{(j+1/2)^2\pi^2+u^2}=\tanh u.
\]
The last equality is the partial-fraction identity obtained by
logarithmically differentiating the Euler product
\[
 \cosh u=\prod_{j\ge0}\left(1+\frac{u^2}{(j+1/2)^2\pi^2}\right).
\]
Termwise integration is justified by
$\int k(\omega)|\sin(\omega u)|\,d\omega<\infty$:
near zero $k(\omega)=O(1/\omega)$, and at infinity it decays
exponentially. The second integral equals
$2\sum_{j\ge0}((j+1/2)\pi)^{-2}=1$.
\end{proof}

\begin{theorem}[A universal work--second-moment transform]
\label{thm:transform}
Suppose random real arithmetic expressions $X(x)$ obey, uniformly on a
finite input set,
\[
 \E X(x)=\mu(x),\quad \E X(x)^2\le V>0,
 \quad \E\operatorname{cost}(X)\le T,\quad T\ge1.
\]
There is a random arithmetic expression $Y(x)$ with mean
$\tanh\mu(x)$ and bounds $V',T'\ge0$, $T'\ge1$, satisfying
\[
 \E Y(x)^2\le V',\qquad
 \E\operatorname{cost}(Y)\le T',\qquad V'T'\le C VT
\]
for a universal constant $C$. The case $V=0$ is handled by the zero
expression. For every $s>0$, the same potential bound holds for the
scaled activation $\tanh(s\mu)/s$.
\end{theorem}
\begin{proof}
For fixed $\omega>0$, put $\lambda=1+\omega^2 V$, take
$K\sim\operatorname{Poisson}(\lambda)$ independently, and define
\[
 Z=\prod_{j=1}^K(1+i\omega X_j/\lambda),\qquad S=\Ima Z.
\]
The empty product is one. The Poisson generating function and fresh
independence give the exact identities
\[
 \E Z=e^{i\omega\mu},\qquad
 \E|Z|^2=\exp\!\left(\frac{\omega^2\E X^2}{\lambda}\right)\le e.
\]
These require only second moments. Since $\mu^2\le V$,
\begin{align*}
 \E S^2&\le\E|Z-1|^2\\
 &=\exp\!\left(\frac{\omega^2\E X^2}{\lambda}\right)
       -2\cos(\omega\mu)+1
 \le C\omega^2 V.
\end{align*}
Together with the preceding bound this proves
$\E S^2\le C\min(1,\omega^2 V)$. The expected work conditional on
$\omega$ is at most $CT\lambda$.

Define
\begin{align*}
 b(\omega)&=\min(1,\omega\sqrt V),\\
 Z_V&=\int_0^\infty k(\omega)b(\omega)/\sqrt\lambda\,d\omega,\\
 A_V&=\int_0^\infty k(\omega)b(\omega)\sqrt\lambda\,d\omega.
\end{align*}
Both integrals are finite and positive. Sample $\omega$ with density
$q_V(\omega)=k(\omega)b(\omega)/(\sqrt\lambda Z_V)$ and output
$Y=k(\omega)S/q_V(\omega)$. The bound on $\E|S|$ gives
absolute integrability, so Fubini and Lemma~\ref{lem:fourier} prove
$\E Y=\tanh\mu$. Moreover
\[
 \E Y^2\le C Z_VA_V,
 \qquad \E\operatorname{cost}(Y)\le CT A_V/Z_V.
\]
For $r\ge0$,
$r\le\min(1,r)\sqrt{1+r^2}\le\sqrt2r$, hence
$\sqrt V\le A_V\le\sqrt{2V}$. Multiplication proves the claim.
Since $A_V/Z_V\ge1$, fixed overhead is absorbed. Apply this
construction to $sX$ and divide its result by $s$ to obtain the last
assertion; the $s^2$ factors cancel in the potential.
\end{proof}

\section{Cost-aware sampling through a network}

\begin{lemma}[Linear combination]
\label{lem:linear}
Suppose child expressions have means $\mu_j$, second-moment bounds
$V_j>0$, and expected-work bounds $T_j\ge1$. A linear combination
$\sum_jw_j\mu_j$ has an unbiased expression with potential at most
\[
 C\left(\sum_j|w_j|\sqrt{V_jT_j}\right)^2.
\]
\end{lemma}
\begin{proof}
Omit zero terms. Set
\[
 L_0=\sum_j|w_j|\sqrt{V_j/T_j},\qquad
 M_0=\sum_j|w_j|\sqrt{V_jT_j},\qquad
 p_j=|w_j|\sqrt{V_j/T_j}/L_0.
\]
Sample $J$ with probabilities $p_j$ and output $w_JX_J/p_J$.
Its mean is the required sum, its second moment is at most $L_0M_0$,
and expected work is at most $CM_0/L_0$. Here $M_0/L_0\ge1$,
which absorbs multiplication overhead. Their product is as claimed.
If all terms vanish, use the zero expression.
\end{proof}

\begin{proposition}[Frobenius propagation]
\label{prop:potential}
A tanh network with $h$ hidden layers, input coordinates of absolute
value at most one, and product $P$ has an unbiased scalar-output
expression whose work--second-moment potential is at most
\[
 C^{h+1}(d+1)P^2.
\]
For a vector readout, the sum of its coordinate potentials has the
same bound. The conclusion also allows neuron-specific normalized
tanh activations or identity activations.
\end{proposition}
\begin{proof}
Input coordinates have $V=T=1$. If $K_j=V_jT_j$ denotes a child
potential, Lemma~\ref{lem:linear} and Theorem~\ref{thm:transform}
give
\[
 K_i\le C\left(\sum_j|w_{ij}|\sqrt{K_j}\right)^2,
 \qquad \sum_iK_i\le C\norm W_F^2\sum_jK_j.
\]
The second inequality is rowwise Cauchy--Schwarz. Iterate it through
the hidden layers and final readout. All multiplication subexpressions
are fresh independent calls; no higher moments are needed.
\end{proof}

The factor $C^{h+1}$ is not a depth-independent constant. The next
normal form is what makes the final estimate polynomial.

\begin{theorem}[Short ordinary-tanh normal form]
\label{thm:normal}
There are universal constants $a,c$ such that a globally unit-margin
tanh Boolean representation of product $P$ has another globally
unit-margin ordinary-tanh representation with at most
$c\log(P+2)$ weighted layers, width, individual weights, and product
at most $(P+2)^a$, using only its at most $P^2$ relevant coordinates.
The weights can be dyadic with at most $c\log^2(P+2)$ bits per entry
after changing the constants.
\end{theorem}

Appendices~\ref{sec:pf-width-compression}--\ref{app:normal} give a
self-contained proof. In particular this statement does not replace
tanh by arbitrary lookup activations. An affine final normal form is
allowed and causes no difficulty in Proposition~\ref{prop:potential}.

\section{Selecting finite formulas and bounding their coefficients}
\label{sec:precision}

Apply Theorem~\ref{thm:normal}. Write $B=d+P+2$ and increase a
universal exponent $a$ as necessary. Its depth is at most $a\log B$;
its dimensions, product, and physical coefficient magnitudes are at
most $B^a$. Its coefficient bit lengths are at most $B^a$. In
particular every nonzero coefficient is at least $\exp(-B^a)$ in
absolute value. The deliberately loose exponential bound is sufficient.

The output estimator $Y$ from Proposition~\ref{prop:potential} has
deterministic bounds $V,T$ with
\begin{equation}\label{eq:VT}
 VT\le B^a.
\end{equation}
The margin implies $V\ge1$, so both final-output bounds are polynomial.
Internal bounds need not individually be polynomial. We next bound
their logarithms, rather than assuming that all internal scales are
well conditioned.

\begin{lemma}[Internal constants]
\label{lem:constants}
All nonzero internal deterministic bounds $V,T,L_0,M_0,Z_V$, and their
reciprocals, lie between $\exp(-B^b)$ and $\exp(B^b)$ for a
universal $b$.
\end{lemma}
\begin{proof}
The explicit estimator bounds may be used as the recursive definitions
of $V$ and $T$, with all work bounds at least one. If preceding
second-moment and work bounds have logarithmic magnitudes at most
$R$, the formulas for $L_0,M_0$ show that their logarithmic
magnitudes are at most $C(R+B^a)$. A nonzero row has a nonzero
coefficient at least $\exp(-B^a)$ and at least one retained child.
Zero children or zero rows are omitted. The same estimate applies to
$V_{\rm in}=L_0M_0$ and $T_{\rm in}=CM_0/L_0$.

The transform uses
\begin{equation}\label{eq:Zbounds}
 \sqrt V\le A_V\le\sqrt{2V},\qquad
 c\min(\sqrt V,1)\le Z_V\le C\min(\sqrt V,1).
\end{equation}
For $V\le1$, integration over $[1,2]$ gives the lower bound for
$Z_V$, and $b(\omega)\le\omega\sqrt V$ gives the upper bound.
For $V\ge1$, integrate over $[0,1/\sqrt V]$ for the lower bound,
using $k(\omega)\asymp1/\omega$ there. For the upper bound split
at $1/\sqrt V$; the remaining integral is at most
$\int_{1/\sqrt V}^\infty k(\omega)/(\omega\sqrt V)\,d\omega$,
which is uniformly bounded by a further split at one. Thus
$V'=CZ_VA_V$ and $T'=CT A_V/Z_V$ satisfy a logarithmic recurrence
$R_{s+1}\le CR_s+CB^a$. There are $O(\log B)$ layers, so its
solution is $B^{O(1)}$.
\end{proof}

\begin{lemma}[A good realization of polynomial size]
\label{lem:realization}
There is a fixed collection of at most $C(d+1)$ arithmetic formulas
of total size polynomial in $B$ such that, on every input, more than
two thirds of their outputs have the correct sign and absolute margin
at least $1/2$. Every coefficient in these formulas has modulus at most
$\exp(B^c)$ for a universal $c$.
\end{lemma}
\begin{proof}
Average $n_0=\lceil40V\rceil$ independent output estimators in
one group. Chebyshev gives error at most $1/2$ with probability at
least $0.9$ at each input. Form $r=O(d+1)$ independent groups, with
a sufficiently large universal constant. Chernoff and a union bound
over $2^d$ inputs give probability at least $0.9$ that more than two
thirds of the groups are good at every input. Their expected total
size is at most $S=B^{O(1)}$, by \eqref{eq:VT} and $V\ge1$.
Markov gives realized size at most $20S$ except with probability
at most $1/20$.

Uniformly over all internal neuron types, Lemma~\ref{lem:constants}
and the density formula imply
\[
 \Prb(\omega<\eta)\le e^{B^b}\eta,\qquad
 \Prb(\omega>U)\le e^{B^b}e^{-c_0U}\quad(U\ge1).
\]
The first estimate uses $k(\omega)\le C/\omega$ and
$b(\omega)\le\omega\sqrt V$; the second uses exponential decay
and the lower bound on $Z_V$. A newly reached node draws its fresh
frequency after its existence and type have been determined. Summing
the conditional bad-frequency probabilities over all tree addresses
therefore bounds the probability of any bad frequency by their common
upper bound times the expected number of reached nodes, at most $S$.
Choose $\eta=e^{-B^{c_1}}$ and $U=B^{c_1}$ with a sufficiently large
universal $c_1$. This failure probability is at most $1/20$.

All three good events have positive joint probability. Fix such a
realization. Every Poisson count is at most its total size. The
arithmetic constants are obtained from original constants and
\[
 \frac{w_j}{p_j}=\sign(w_j)L_0\sqrt{T_j/V_j},\quad
 \frac\omega\lambda,\quad
 \frac{k(\omega)}{q_V(\omega)}
       =\frac{Z_V\sqrt\lambda}{\min(1,\omega\sqrt V)},\quad
 n_0^{-1}.
\]
Their logarithmic magnitudes are polynomial in $B$ on the good event.
Each good group approximates an output of absolute value at least one
to error at most $1/2$, giving the margin. We have selected an event
from the original unbiased construction; we have not conditioned the
estimator on truncated frequencies and then asserted unbiasedness.
\end{proof}

\begin{lemma}[Dyadic rounding and integer numerators]
\label{lem:round}
The formulas in Lemma~\ref{lem:realization} can be replaced by
polynomial-size formulas over Gaussian integer constants, addition,
multiplication, and real/imaginary-part gates, with all integer constants
of polynomial bit length, so that their majority of signs still equals
$f$ on the whole cube.
\end{lemma}
\begin{proof}
Regard each occurrence of a constant in a formula as a separate pair
of real variables. A formula of size $s$ is separately linear in each
occurrence, with degree at most $s$ and at most exponentially many
terms. Real/imaginary-part gates are real-linear. On a coefficient box
of magnitude $e^{B^{O(1)}}$, values and first coefficient derivatives
are therefore at most $e^{B^{O(1)}}$ because $s=B^{O(1)}$.
Equivalently this follows by direct induction through the gates.
Round every coefficient component to dyadic error $2^{-B^c}$ with
a sufficiently large constant exponent. Every group output changes by
less than $1/8$, uniformly on the cube. Its previously good outputs
retain the correct sign and margin at least $3/8$. Bad groups may
change sign, but more than two thirds remain good at every input.

Clear denominators with positive integer subtree denominators. At an
addition replace $N_1/D_1+N_2/D_2$ by
$(N_1D_2+N_2D_1)/(D_1D_2)$; at multiplication multiply numerators
and denominators. Real and imaginary part commute with division by
a positive real denominator. Each denominator is a precomputed
constant whose bit length is at most the sum of leaf denominator bit
lengths in the subtree. Each child numerator occurs only once, so the
result is a polynomial-size formula, not an exponentially expanded
circuit. The sign of its real output is that of its integer numerator.
\end{proof}

\section{Balancing and the arithmetic complexity class}

\begin{lemma}[Real-affine balancing of complex formulas]
\label{lem:balancing}
A size-$s$ formula over Gaussian integer constants, complex addition
and multiplication, and unary real/imaginary-part gates has an
equivalent real integer formula of polynomial size and depth $O(\log(s+2))$.
Its constants have no larger bit lengths, up to a constant, than the
original constants.
\end{lemma}
\begin{proof}
First balance the complex formula. Select a subtree $Q$ whose size
lies between one third and two thirds of the full size, up to a fixed
additive constant. Such a subtree exists for a tree of fan-in at most
two, including unary gates. Since $Q$ occurs only once, its surrounding
context is real-affine. Thus
\[
 F=(F[Q:=1]-F[Q:=0])\Rea Q
   +(F[Q:=i]-F[Q:=0])\Ima Q+F[Q:=0].
\]
Both $Q$ and each substituted context have at most two thirds of
the original size, up to fixed additive terms. Recursing introduces
only a fixed number of copies at each recursive level. The resulting
size is polynomial in $s$ and depth is $O(\log(s+2))$. The only new
constants are $0,1,i,-1$.

Now replace each complex gate by its fixed-size real/imaginary pair
of integer arithmetic gates. Expansion into formulas duplicates
subexpressions by only a fixed factor per level. Since the balanced
depth is logarithmic, this gives only a polynomial total factor.
Doing this expansion before balancing would not justify the bound.
The final output is real, so keep its real component.
\end{proof}

\begin{proof}[Proof of Theorem~\ref{thm:main}]
Lemmas~\ref{lem:realization}--\ref{lem:balancing} give polynomially
many polynomial-size logarithmic-depth integer formulas. Their
polynomial-bit constants can be built using $-1,0,1$: construct a
power $2^j$ as a balanced product of $j$ copies of $1+1$, and express
an integer as a balanced sum of its signed nonzero powers of two.
This adds polynomial size and logarithmic depth. Replace sign inputs
by $2b-1$ for bit inputs. Each formula's positivity predicate belongs
to nonuniform $\PNC^1$.

The final majority is itself a $\PNC^1$ predicate (indeed it is the
sign of an affine integer sum of the input bits). It is applied to
polynomially many $\PNC^1$ predicates. The composition theorem
$[\PNC^1]\cdot[\PNC^1]=\PNC^1$ in \cite[Theorem 4.2]{datta}
therefore gives the claimed circuit with universal polynomial size
and logarithmic depth. Equivalently use the constant nesting depth
of positivity gates and \cite[Proposition 3.4 and Theorem 4.1]{datta}.
The inclusion in $\Lpoly$ follows from their Proposition 2.4, with
the selected circuit supplied as advice. Formula balancing alone
would not establish Boolean $\NC^1$.
If the final integer output can be zero, replacing it by twice that
output minus one gives strictly positive and strictly negative values
on the two label classes without changing the resource bounds.
\end{proof}

\section{A converse for the original Gaussian posterior}
\label{sec:posterior}

Fix $d,n\ge1$, learner weighted depth $L\ge3$, and common hidden
width $N$. Write $\sigma^2$ for the common variance in energy
coordinates. The independent physical weight variances are
\[
 W^1_{ij}:\frac{N\sigma^2}{d},\qquad
 W^s_{ij}:\sigma^2\ (2\le s<L),\qquad
 v_i:\frac{\sigma^2}{N}.
\]
Equivalently the Gaussian energy is
\[
 \kappa_N(W)=\frac dN\norm{W^1}_F^2+
       \sum_{s=2}^{L-1}\norm{W^s}_F^2+N\norm v_2^2.
\]
For a consistently labeled sample $S=((x_i,y_i))$, the posterior
$Q_{S,\sigma}$ conditions the raw prior on
\[
 |z(x)|>1/2\quad\hbox{for every }x\in X_d,
 \qquad y_i z(x_i)>0\quad\hbox{for every training example}.
\]
No hidden robustness, internal margin, or parameter bound is imposed.
Its induced distribution on Boolean truth tables will also be denoted
by $Q_{S,\sigma}$. Define
$p_{S,\sigma}(f(x)\mid x)=Q_{S,\sigma}\{h:h(x)=f(x)\}$.
The exponent convention is $\sigma^2=N^{-\gamma}$.

\subsection{The precise imported posterior input}

Let $\mathcal H_P$ be the finite set of Boolean functions on $X_d$
admitting some globally unit-margin tanh representation of product at
most $P$, with arbitrary finite depth and widths. The following
posterior theorem is an input from the existing Gaussian analysis
\cite{unrestricted}; it is not a consequence of the Fourier estimator.
We state all its quantitative data so its use here is unambiguous.

\begin{theorem}[Finite-scale representation ceiling, imported]
\label{thm:ceiling}
Fix $\zeta>0$ and $0<\beta<1$. Define
\begin{align*}
 C_0&=n+1,\qquad R_0=\max\{4(n+1),(d+1)\log(16(n+1))\},\\
 M&=\left\lceil\max\{4,C_0,2R_0,dC_0(d+1)R_0^2\}\right\rceil,
 \qquad T_0=MC_0R_0^2,\\
 K_0&=2+(L+14)(2T_0)^{1/(L-1)},\\
 b_0^2&=\max\{2(d+1)\log2+4\log8,\ 16\log9+8\log(8L)\},\\
 A_0&=\left(\frac{K_0+2}{L-1}\right)^{(L-1)/2},\\
 \xi&=\min\left\{1,\frac1{4\sqrt{L(K_0+1)}},
                     \frac1{8\sqrt2 LA_0}\right\},\\
 N_0&=\left\lceil\max\{M,2,(b_0/\xi)^{2/\zeta},
                                     2\log(2/\beta)\}\right\rceil,\\
 P_0^2&=d^{-1}e^{19}(2T_0)^{3/2}.
\end{align*}
For every $N\ge N_0$, every actual positive scale
$\sigma^2\le N^{-(1+\zeta)}$, and every consistently labeled
sample of at most $n$ inputs,
\[
 Q_{S,\sigma}(\mathcal H_{P_0})\ge1-\beta.
\]
The threshold $N_0$ is polynomial in
$d,n,L,\log(2/\beta)$ with degree depending only on $\zeta$.
The product bound satisfies
\[
 P_0\le e^{19/2}544^{3/4}(d+n+2)^9,
\]
independently of $L,N,\zeta,\beta$.
\end{theorem}

The low-product representative can differ from the sampled parameter
vector. The theorem does not assert that posterior weights themselves
have bounded energy. Its proof uses a uniform memorizing witness,
a Gaussian perturbation set, and Gaussian distance concentration.
The exact source is the theorem labeled
\texttt{thm:posterior-energy-ceiling} in the file identified in
\cite{unrestricted}; the package dependency manifest records its hash.
For the final displayed polynomial, put $S_0=d+n+2$. Directly,
$C_0\le S_0$, $R_0\le4S_0^2$, $M\le17S_0^7$, and
$T_0\le272S_0^{12}$, which give the claim.

\subsection{Two majority reductions}

\begin{theorem}[From distribution-free samples to arithmetic circuits]
\label{thm:converse}
Fix $d,n\ge1$, $L\ge3$, $\zeta>0$, and $0\le\varepsilon<1/2$.
Use Theorem~\ref{thm:ceiling} with $\beta=1/6$, and choose one
$N\ge N_0$ and one $0<\sigma^2\le N^{-(1+\zeta)}$, before
choosing an input distribution. Suppose that for every probability
distribution $\mathcal D$ on $X_d$, an iid $n$-sample labeled by $f$
has positive probability of satisfying
\begin{equation}\label{eq:confidence}
 \mathcal D\{x:p_{S,\sigma}(f(x)\mid x)\le2/3\}\le\varepsilon.
\end{equation}
Then $f$ has an integer arithmetic circuit with constants $-1,0,1$
whose sign equals $f$, size polynomial in
$d+n+(1-2\varepsilon)^{-1}$, and depth
$O(\log(2+d+n+(1-2\varepsilon)^{-1}))$.
All polynomial exponents are universal and independent of learner
depth, width, scale, and $\zeta$.
\end{theorem}
\begin{proof}
Fix $\mathcal D$ and choose a successful sample. Condition its
posterior on $\mathcal H_{P_0}$, obtaining $\nu$. On
$G=\{x:p_{S,\sigma}(f(x)\mid x)>2/3\}$, which has
$\mathcal D$-mass at least $1-\varepsilon$, its wrong-label
probability is strictly below $(1/3)/(5/6)=2/5$. Thus
$\E_{h\sim\nu}f(x)h(x)>1/5$ at every $x\in G$.

Choose an odd $m_0$ between $100(d+1)$ and $100(d+1)+2$.
Hoeffding bounds the error of $m_0$ independent voters at a fixed
good input by $e^{-m_0/50}$. The union bound over $2^d$ inputs is
less than one. Some fixed majority $g$ of these low-product functions
therefore agrees with $f$ everywhere on $G$ and has
$\mathcal D$-error at most $\varepsilon$.

Let $\mathcal K$ be the finite class of all $m_0$-voter majorities
of $\mathcal H_{P_0}$. It does not depend on $\mathcal D$.
Writing $\kappa=1-2\varepsilon>0$, we have proved
\[
 \text{for every }\mathcal D,\qquad
 \max_{g\in\mathcal K}\E_{x\sim\mathcal D}f(x)g(x)\ge\kappa.
\]
Finite minimax gives a distribution $\lambda$ on $\mathcal K$
with $f(x)\E_\lambda g(x)\ge\kappa$ for every input. Choose
an odd $m_1\ge4(d+1)/\kappa^2$ within two of that value.
Another Hoeffding bound and union bound give
$2^d e^{-m_1\kappa^2/2}<1$. Some majority of $m_1$ members of
$\mathcal K$ is consequently equal to $f$ on the entire cube.

Compile its $m_0m_1$ low-product constituents by
Theorem~\ref{thm:main}. Their resources are polynomial in $d+n$
by Theorem~\ref{thm:ceiling}. Two further majority layers are a
constant number of $\PNC^1$ predicate-composition levels, absorbed
by the closure theorem \cite{datta}. The asserted polynomial and
depth bounds follow. No assertion that a majority itself has low
Frobenius product has been used.
\end{proof}

The proof handles posterior mixtures without choosing a unique
minimum-energy function. The first majority uses confidence on a large
set; the second uses distribution-universal approximation and minimax.
Averaging the error of individual posterior draws directly would give
a weaker range of $\varepsilon$ and would miss this distinction.
Only existence of a successful sample is used, so any usual
positive-probability or high-probability sample guarantee implies the
hypothesis. Its distribution-free quantifier is essential.

\subsection{Fixed exponents and the strong-prior limit}

\begin{corollary}[The two strong-prior regimes]
\label{cor:regimes}
The conclusion of Theorem~\ref{thm:converse} holds in each case below.
\begin{enumerate}[label=(\roman*)]
\item Fix $\gamma>1$, put $\zeta=\gamma-1$, choose any
$N\ge N_0$, and use $\sigma^2=N^{-\gamma}$.
\item Fix $d,n,L$, choose $N\ge N_0$ using $\zeta=1$ and
$\beta=1/6$, and let $Q_S^0$ be any selected strong-prior limit of
the posterior truth-table distributions as $\sigma\downarrow0$.
Assume \eqref{eq:confidence} for $Q_S^0$ directly.
\end{enumerate}
In (ii), a common subsequence for all samples may be used, and the
claim includes an actual posterior limit whenever it exists.
\end{corollary}
\begin{proof}
Case (i) is immediate. For (ii), every sufficiently small $\sigma$
satisfies Theorem~\ref{thm:ceiling}. Truth-table space is finite and
$\mathcal H_{P_0}$ is a fixed subset, so any limiting posterior has
$Q_S^0(\mathcal H_{P_0})\ge5/6$. There are only finitely many
samples of at most $n$ inputs, so sequential compactness also permits
a common subsequence. Run the proof of Theorem~\ref{thm:converse}
on these limiting probability vectors. No passage of the strict
confidence event through a limit, or independent assertion of limit
uniqueness, is needed.
\end{proof}

The strong-prior width threshold above is sufficient; it is not
claimed optimal. In both cases the chosen architecture precedes
the input distribution. A proposed sample bound valid at every
realizable width permits choosing this sufficiently large width after
its width-independent sample budget has been fixed.

\begin{corollary}[A conditional barrier to universal polynomial learning]
\label{cor:collapse}
Fix either a Gaussian exponent $\gamma>1$ or the strong-prior limit
of Corollary~\ref{cor:regimes}. Suppose the original learner has a
distribution-free sample guarantee with some fixed
$\varepsilon<1/2$ and positive sample-success probability, bounded
by one polynomial jointly in input dimension, bounded-weight teacher
width and depth, and learner depth, independently of learner width,
valid at every realizable learner width. Then
\[
 \Ppoly=\text{nonuniform }\PNC^1=\Lpoly.
\]
\end{corollary}
\begin{proof}
Every polynomial-size Boolean circuit family has a tanh teacher
family with polynomial width and depth and uniformly bounded physical
weights. Here is an explicit elementary construction. Compile gates
to NAND. Carry each already computed logical bit by a neuron with
activation within $1/10$ of its sign, and carry a constant $1/2$
using the weight $2\operatorname{atanh}(1/2)$ after its first layer.
The map $\tanh(4u)$ copies a bit. A NAND of $u,v$ is
$\tanh(4(1-u-v))$, with the affine constant implemented by the
charged constant wire. Its preactivation has the required sign and
absolute value at least $3.2$, so its output is again within $1/10$
of the desired sign. Initialize inputs by $\tanh(4x_i)$, execute
one gate per layer while carrying previous bits, and multiply the
final bit by two. Width and depth are polynomial in the Boolean
circuit size, all weights have magnitude at most eight, and the
global output margin exceeds one. Zero padding to larger widths
preserves the function. If a convention requires teacher weights
bounded by teacher width, pad the width to at least eight.

The hypothesized sample budget is therefore polynomial in $d$ for
every family in $\Ppoly$. Choose an admissible learner width above
the ceiling threshold and apply Theorem~\ref{thm:converse} or
Corollary~\ref{cor:regimes}. This puts the family in nonuniform
$\PNC^1\subseteq\Lpoly$. The reverse inclusions in $\Ppoly$
are standard finite-circuit simulations, proving equality.
\end{proof}

This is not a known contradiction and not an exponential sample lower
bound. Even a conditional failure of every polynomial bound is weaker
than an $\exp(\Omega(L))$ lower bound in teacher depth. A sample
guarantee polynomial in the chosen learner width, or a guarantee only
for one input law, would not imply the displayed collapse.

\begin{remark}[Other confidence thresholds]
For a fixed threshold $\tau>1/2$, choose
$\beta=\tau-1/2$ in the ceiling. Conditioning then leaves signed
vote gap at least
$(2\tau-1-\beta)/(1-\beta)>0$. The first majority has size
$O((d+1)/\mathrm{gap}^2)$, and the same argument applies. At
$\tau=1/2$ there is no uniform positive vote gap supplied by this
calculation; the same polynomial conclusion is not proved here.
\end{remark}


\section{Scope and open questions}

The result isolates a tanh-specific structural restriction on the
functions favored by sufficiently strong priors. Its proof is
nonuniform and allows enormous original weights and widths, yet the
replacement arithmetic circuit has polynomially many bits. It is not
an algorithm for posterior sampling or circuit recovery.

The endpoint $\gamma=1$ remains outside the new posterior converse.
Existing descriptions of its low-action functions can use nonlinear
expectation maps. A small value of that action is not a proved
low-product deterministic tanh representation, so applying
Theorem~\ref{thm:main} to it would be an extra, unsupported step.
Related circuit upper bounds in larger classes do not fill that gap.

Nor does the theorem prove that some small Boolean circuit needs
exponential product or samples. It gives a conditional obstruction:
universal polynomial learning would impose a strong nonuniform
complexity collapse. It leaves open intermediate growth rates and
does not identify an explicit hard target. A converse inclusion such
as $\PNC^1$ inside the low-product tanh class is also not established.
Finally, the scalar transform extends to normalized tanh, but its
Fourier formula does not automatically cover ReLU, GELU, or arbitrary
Lipschitz activations.

\appendix
\section{Width compression using ordinary tanh}
\label{sec:pf-width-compression}

We give the finite-set compression argument in full because preserving
ordinary tanh, rather than merely preserving Lipschitz activations, is
essential for the Fourier estimator. The proof also applies to only the
first $h$ hidden layers: replace $H$ by $h$ in its error and norm budgets
and stop after those $h$ steps. None of the resulting prefix-width bounds
depends on the number or widths of later layers.

\begin{theorem}[Finite-set width compression for tanh]
\label{thm:pf-width-compression}
Let $X$ be a nonempty finite set of inputs of norm at most $R\ge1$.
Suppose a depth-$D_0\ge2$ tanh network with linear readout, arbitrary
finite hidden widths, and no later biases has product of Frobenius norms
at most $P_0\ge1$. For every $0<\epsilon<1$, it has a same-depth
approximation $\widetilde z$ such that
\[
 \max_{x\in X}|z(x)-\widetilde z(x)|\le\epsilon,
 \qquad \prod_{j=1}^{D_0}\|\widetilde W^j\|_F\le2P_0,
\]
and every new hidden width is bounded by a universal polynomial in
\[
 P_0,D_0,R,\log(2|X|),\epsilon^{-1}.
\]
The bound is independent of the original widths and individual weights.
An augmented constant input is allowed and counts toward $R$.
\end{theorem}

\begin{proof}
Put $H=D_0-1$, $P_*=2P_0$, $e=\epsilon/(8H)$ and
$\eta=1/(4H)$. We compress the hidden layers in order from the first
to the last. At each step the product norm increases by at most
$1+\eta$, and the output changes by at most $6e$ on $X$.
Thus every intermediate product is below $P_*$, and the final product
is at most $e^{1/4}P_0<2P_0$, while total error is less than $\epsilon$.
A zero matrix makes the entire network zero and is trivial, so all
prefix and pair norms used below may be assumed positive. If a
compression step produces a zero matrix, stop and use the zero network
with width one; the accumulated approximation estimates still apply.

Consider a hidden matrix $W$ and its immediately following matrix $V$;
the latter may be the final linear readout. Let $S$ be the product of
all Frobenius norms strictly before $W$, with $S=1$ at the first layer,
and put
\[
 A=\|W\|_F\|V\|_F.
\]
The incoming feature vector $h(x)$ has norm at most $SR$. The part
strictly after $V$ has Lipschitz constant at most $P_*/(SA)$.
If $w_i^T$ is row $i$ of $W$ and $v_i$ is column $i$ of $V$, write
$a_i=\|w_i\|_2$ and $b_i=\|v_i\|_2$.
Zero rows or zero outgoing columns can be deleted.

Set
\[
 B=P_*/e,\qquad \tau^2=3e/(P_*R^3),\qquad
 \Delta=e/(P_*R),\qquad Q=(B/\tau)^2.
\]
Partition neurons into large, small and middle according as
$Sa_i>B$, $Sa_i<\tau$, or $\tau\le Sa_i\le B$.
There are at most $S^2\|W\|_F^2/B^2$ large neurons. Their activation
vector has norm at most the square root of this count. Deleting them
therefore changes the final output by at most $P_*/B=e$.

For the small neurons, replace tanh by the identity. Since
$|u-\tanh u|\le |u|^3/3$, their combined activation error has norm at
most $\tau^2R^3S\|W\|_F/3$. Propagation through the suffix bounds
the output error by $P_*\tau^2R^3/3=e$.
Their resulting contribution to the next preactivation is the linear map
\[
 B_0=\sum_{i\ \mathrm{small}}v_iw_i^T,
 \qquad \|B_0\|_*\le\sum_{i\ \mathrm{small}}a_ib_i,
\]
where $\|\cdot\|_*$ is the nuclear norm.

Let $A_m=\sum_{i\ \mathrm{middle}}a_ib_i$.
If $A_m<\Delta A$, delete all middle neurons: their output contribution
is at most $P_*R A_m/A\le e$. Otherwise $A_m\ge\Delta A>0$, and
we approximate their contribution by sampling. Define
\[
 Z=\sum_{i\ \mathrm{middle}}b_i/a_i,
 \qquad p_i=(b_i/a_i)/Z.
\]
Draw $m$ independent indices from $(p_i)$. A sampled neuron keeps its
incoming row $w_i^T$ and has outgoing column $v_i/(mp_i)$.
If $\widehat A=(Z/m)\sum_{r=1}^m a_{i_r}^2$ and $c=m/Z$, its new
incoming and outgoing squared Frobenius norms are exactly
\[
 c\widehat A\quad\hbox{and}\quad \widehat A/c.
\]
Moreover $\mathbb E\widehat A=A_m$. Since middle row norms lie between
$\tau/S$ and $B/S$, one has
\[
 Z\max_i a_i^2\le A_m Q\le AQ.
\]
The random vector $v_i\tanh(w_i^Th(x))/p_i$ has norm at most
$AQSR$, uniformly on $X$.

Consequently some choice of the samples has both
\[
 \widehat A\le A_m+\eta A
\]
and middle-preactivation approximation error at most $eSA/P_*$
simultaneously for every $x\in X$, provided
\[
 m\ge64Q^2\max\{\eta^{-2},(P_*R/e)^2\}
                       \log(8(|X|+1)).
\]
Here and below take integer ceilings. For detail, the scalar random
variables $Za_i^2/A$ are bounded by $Q$, so Hoeffding controls their
mean. For a Hilbert-valued random vector bounded by $AQSR$, the mean
deviation has expectation at most $AQSR/\sqrt m$; bounded differences
gives a dimension-free tail. A union bound over $X$ and the scalar
condition has positive success probability with the displayed constant.
Propagation of the vector error gives final output error at most $e$.

We next compress the linear branch. Truncate the singular-value
decomposition of $B_0$ to rank
\[
 k=\lceil P_*R/e\rceil,
\]
obtaining $B_k=UV_0$ with
$\|U\|_F^2=\|V_0\|_F^2=\nu:=\|B_k\|_*\le\|B_0\|_*$.
The residual operator norm is at most $\|B_0\|_* /(k+1)\le A/(k+1)$,
so its final output error is at most $P_*R/(k+1)\le e$.
If $\nu=0$, omit the linear branch entirely.

When the middle branch was sampled, retain $c=m/Z$. When it was
deleted and $\nu>0$, choose $c=1/(\nu S^2)$ instead.
Represent $B_k$ approximately using $K$ copies of each of its at most
$k$ singular channels: their incoming rows are the rows of
$\sqrt{c/K}\,V_0$, repeated $K$ times, and their outgoing columns are
the corresponding columns of $U/\sqrt{cK}$.
Their total incoming and outgoing squared norms are $c\nu$ and
$\nu/c$. Their linear limit is exactly $B_kh$.
The tanh linearization error at the next preactivation has norm at most
\[
 \frac{c\nu^2(SR)^3}{3K}.
\]
After propagation this is at most
$P_*c\nu S^2R^3/(3K)$, using $\nu\le A$.
In the sampled case,
\[
 Z\ge A_mS^2/B^2,
 \qquad c\nu S^2\le mB^2/\Delta;
\]
in the deleted case this quantity equals one. Hence it suffices to use
\[
 K=\left\lceil\frac{P_*R^3}{e}
                   \max\{1,mB^2/\Delta\}\right\rceil.
\]
The linear-channel approximation then contributes at most $e$ error.
When there is no sampled middle branch, the same bound is harmless;
one may use the displayed deterministic value of $m$ merely to define
a common width upper bound.

Combining the sampled and linear branches gives pair product norm
\[
 \widehat A+\nu
 \le A_m+\sum_{i\ \mathrm{small}}a_ib_i+\eta A
 \le(1+\eta)A.
\]
The last inequality is Cauchy--Schwarz over the original rows and
columns. If the middle branch was deleted, the pair product is simply
$\nu\le A$. These are exact norm calculations for the constructed
tanh matrices, not only bounds on their linear limits.

The compressed hidden width is at most $m+kK$. All quantities in this
bound are universal polynomials in $P_0,H,R,e^{-1}$ and
$\log(2|X|)$. Crucially, no original dimension of either adjacent hidden
layer occurs. Only the two adjacent matrices are altered. Processing
the next hidden layer changes its own width without undoing the width
bound for the preceding one. Uniform zero-padding at the end makes
all hidden widths equal if desired. Summing the at most six errors and
multiplying the pair-product bounds proves the theorem.
\end{proof}

\section{A logarithmic-depth ordinary-tanh normal form}
\label{app:normal}

We supply the geometric reduction and all remaining normalization
steps. They are needed because the estimator acts on tanh itself,
whereas a generic finite-precision simulation could replace tanh by
lookup tables without preserving the hypothesis of the estimator.

\subsection{Relevant coordinates and a two-ball cover}

\begin{lemma}[Relevant-coordinate reduction]
\label{lem:relevant}
A globally unit-margin tanh representation has $P\ge1$, and its
Boolean function has at most $P^2$ relevant input coordinates.
Fixing all the other coordinates can be done without increasing its
product or reducing its global margin on the remaining cube.
\end{lemma}
\begin{proof}
Since tanh is zero-preserving and one-Lipschitz,
\[
 1\le\E_x z(x)^2\le
 \left(\prod_{s=2}^D\norm{W^s}_F^2\right)
 \E_x\norm{W^1(1,x)}_2^2=P^2,
\]
where the expectation is uniform and
$\E[(1,x)(1,x)^T]=I$. Put $T=\prod_{s=2}^D\norm{W^s}_F$.
For every relevant coordinate $i$, a bichromatic cube edge in that
direction gives $2\le2T\norm{W^1e_i}_2$. Summing squares yields
at most $P^2$ relevant coordinates. For the irrelevant columns,
choose independent uniform signs $a_i$, replace the first constant
column $w_0$ by $w_0+\sum_i a_iw_i$, and delete those columns.
The expected new squared Frobenius norm equals the old squared norm.
Some choice does not increase it. Every restriction has the same
Boolean labels and retains the original margin.
\end{proof}

\begin{lemma}[Two balls]
\label{lem:twoball}
For every real matrix $W$, every Euclidean ball $B(a,r)$, and every
coordinatewise one-Lipschitz map $\Phi$, the set
$\Phi(W B(a,r))$ is covered by two balls of radius
$r\norm W_F/\sqrt2$.
\end{lemma}
\begin{proof}
Translation and scaling reduce to $a=0,r=1,\norm W_F=1$;
the zero case is immediate. Let $\alpha$ be the leading singular
value. If $\alpha^2\le1/2$, one ball suffices. Otherwise put
$b^2=1-\alpha^2$ and $c=(2\alpha^2-1)/(2\alpha)$.
The remaining singular part has operator norm at most $b$.
With $u$ a leading left singular vector, a point whose leading
input coordinate is $t$ has squared distance to the matching
center $cu$ or $-cu$ at most
\[
 (\alpha|t|-c)^2+b^2(1-t^2).
\]
This is convex in $|t|\in[0,1]$ and both endpoint values equal
$1/(4\alpha^2)\le1/2$. Apply $\Phi$ to the two centers;
its Euclidean Lipschitz constant is at most one.
\end{proof}

For prefix products $S_j=\prod_{s\le j}\norm{W^s}_F$ and
inputs of norm at most $R$, iteration gives at most $2^j$ balls of
radius $RS_j2^{-j/2}$. The remaining suffix is Lipschitz with
constant at most $P/S_j$. Thus normalized prefix states of opposite
labels are at least $2/P$ apart. If
\[
 j=\left\lceil2\log_2(64PR)\right\rceil,
\]
selecting one actual state from each nonempty ball gives polynomially
many centers $c_a=h^j(x_a)/S_j$, each with its label $f_a$.
Every normalized state has a same-label center within $1/(8P)$,
and every opposite-label center is at least $2/P$ away. All states
and centers have norm at most $R$. The center need not initially lie
in the image: selecting an actual image doubles the cover radius,
which the displayed constants allow. Every covering ball is
monochromatic by the suffix Lipschitz estimate.

\subsection{Physical weights and finite precision at short depth}

\begin{lemma}[Bounding physical weights]
\label{lem:physical}
Suppose a globally unit-margin tanh representation already has
depth $D$, hidden widths at most $C$, and product at most $P$.
It has a representation of the same Boolean function with no more
weighted layers, product at most $2P$, margin greater than $3/2$,
and individual weights bounded by a universal polynomial in
$d,D,C,P$. It can then be rounded to dyadic weights with
$O(D\log B)$ bits and product at most $4P$, where $B$ is a
polynomial bounding its dimensions, Frobenius norms, and row sums.
The rounded representation retains unit margin.
\end{lemma}
\begin{proof}
Write $M_j=\norm{W^j}_F$, $S_j=\prod_{s\le j}M_s$, and
$R=\sqrt{d+1}$. Margin and bounded hidden activations give
$S_j\le P\sqrt C$ at every hidden layer. Set
\[
 \tau^2=\frac1{16PD R^3}.
\]
Replace tanh by the identity at each hidden layer whose original
prefix product is less than $\tau$. For any intermediate hybrid
network, all prefix and suffix Lipschitz bounds remain valid.
Using $|u-\tanh u|\le|u|^3/3$ and
$\norm{v^{\odot3}}_2\le\norm v_2^3$, the telescoping output
error is at most
\[
 \sum_{j:S_j<\tau}\frac{P}{S_j}\frac{(S_jR)^3}{3}
 \le PD\tau^2R^3/3\le1/48.
\]
Collapse consecutive linear blocks into the following nonlinear
layer or final readout, then double the readout. If retained hidden
indices are $j_1<\cdots<j_H$, the collapsed matrices have assigned
Frobenius bounds
$S_{j_1}$, $S_{j_s}/S_{j_{s-1}}$, and $2P/S_{j_H}$.
Since each retained $S_{j_s}$ lies between $\tau$ and $P\sqrt C$,
all weights are bounded by $\max(1,2P\sqrt C/\tau)$, a polynomial.
The product is at most $2P$ and the margin exceeds $3/2$.
If no hidden activation remains, the affine matrix has norm at
most $2P$ and the same conclusions hold.

For precision, enlarge a polynomial $B\ge2$ to bound all dimensions,
matrix Frobenius norms, and absolute row sums. At remaining depth
$D'\le D$, every nonzero matrix norm satisfies
\[
 M_j\ge1/(\sqrt{d+1}B^{D'-1}),
\]
because $1\le|z(x)|\le\sqrt{d+1}\prod_jM_j$.
Round entries to mesh $2^{-b}$ with
$b=O(D'\log B+\log(d+D'+2))$, with a sufficiently large
universal constant. Matrix perturbations are at most $M_j/(4D')$,
so the product increases by less than two. The layerwise error
recursion using the row-sum bound $B$ and Lipschitzness of tanh
makes the output error less than $1/2$ at the same precision.
The pre-rounding margin in fact exceeds $2(1-1/48)>3/2$, so the
remaining margin is greater than one. If needed, absorb the additive
logarithm into $O(D\log B)$ by enlarging $B$.
\end{proof}

\subsection{Replacing the long suffix by a tanh classifier}

\begin{proof}[Proof of Theorem~\ref{thm:normal}]
Apply Lemma~\ref{lem:relevant}, retain $k\le P^2$ inputs, and
put $R=\sqrt{k+1}$ and
$j=\lceil2\log_2(256PR)\rceil$.
If $D=1$, the affine coefficients have norm at most $P$; double
the readout and round them directly to sufficiently fine dyadic
precision $O(\log(P+2))$, retaining unit margin, and finish.
Hence assume $D\ge2$.
Use the prefix version of the compression theorem in
Appendix~\ref{sec:pf-width-compression} on the first
$\min(j,D-1)$ hidden layers with error $1/4$, then double the
readout. The global margin is at least $3/2$, the full product is
at most $P_0=4P$, and all retained widths are bounded by a
universal polynomial $C=C(P)$: the only inputs to that width
bound are $P,R,j,\log(2|X_k|)$ and a fixed error tolerance.
If all hidden layers have been retained, apply
Lemma~\ref{lem:physical} and finish. Otherwise the first $j$
hidden layers have the cover and normalized centers just proved,
with $P_0$ in place of $P$.

First supply a constant to a new suffix while preserving the product
bound. For retained prefix products $S_s$, global margin gives
$S_s\le P_0\sqrt C$. Choose
\[
 \epsilon_0=1/\max(1,2\max_{s\le j}S_s).
\]
It has inverse polynomial in $P$ and satisfies
$0<\epsilon_0S_s\le1/2$. Add one constant rail whose activation
in layer $s$ is exactly $\epsilon_0S_s$. Its first weight is
$\operatorname{atanh}(\epsilon_0S_1)$ on the augmented constant;
its later weight is
$\operatorname{atanh}(\epsilon_0S_s)/(\epsilon_0S_{s-1})$.
For $0<u\le1/2$,
$(\operatorname{atanh}u/u)^2\le1+u^2\le5/4$.
Thus the added squared norm is at most $(5/4)M_s^2$ in each
layer, and the augmented prefix product is at most $(3/2)^jS_j$.

At the interface the features are $(h^j,\epsilon_0S_j)$.
A suffix written for $(u,1)=(h^j/S_j,1)$ is attached by multiplying
its first matrix by
$\operatorname{diag}(S_j^{-1}I,(\epsilon_0S_j)^{-1})$.
This increases its first Frobenius norm by at most
$(\epsilon_0S_j)^{-1}$. The complete product is therefore at
most $(3/2)^j\epsilon_0^{-1}$ times the suffix's ordinary product.
Both extra factors are polynomial in $P$.

We now give the suffix. Its input $u$ and each of the $q=\poly(P)$
centers $c_a$ have norm at most $R$. A correct-label center is
within $1/(8P_0)$, while each opposite-label center is at distance
at least $2/P_0$. In one hidden layer provide the features
\[
 \tanh(\beta u_i),\qquad
 \tanh(1+\alpha u_i),\qquad\tanh(1-\alpha u_i),
\]
and an exact constant $1/2$. Linear combinations of these features
approximate respectively $u_i$ and $u_i^2$ using
\[
 \frac{\tanh(\beta u_i)}\beta,\qquad
 \frac{\tanh(1+\alpha u_i)+\tanh(1-\alpha u_i)-2\tanh1}
      {\alpha^2\tanh''1}.
\]
Their errors are $O(\beta^2R^3)$ and $O(\alpha^2R^4)$.
Choose $\alpha,\beta$ inverse-polynomially small in
$C,R,P_0$ so that every affine combination
\[
 D_a(u)=\sum_i\operatorname{approx}(u_i^2)
       -2\sum_i(c_a)_i\operatorname{approx}(u_i)+\norm{c_a}_2^2
\]
approximates $\norm{u-c_a}_2^2$ to error at most
$1/(256P_0^2)$. This is possible with polynomial weights because
$\norm{c_a}_1\le\sqrt C R$ and $\tanh''1\ne0$.
Every affine constant is charged to the constant rail.

In a second hidden layer use
\[
 t_a(u)=\tanh(B_0[\tau_0-D_a(u)]),\qquad
 \tau_0=1/(16P_0^2),\quad B_0=32P_0^2\log(16q),
\]
and propagate the constant $1/2$. Put $\chi_a=(1+t_a)/2$.
For a selected close center, $\tau_0-D_a\ge11/(256P_0^2)$;
for every opposite-label center it is less than $-3/P_0^2$.
The tanh tail estimate gives $\chi_a\ge1-1/(16q)$ for a close
center and $\chi_a\le1/(16q)$ for each opposite-label center.
All other $\chi_a$ lie in $[0,1]$, so
\[
 f(u)\sum_a f_a\chi_a(u)\ge1-1/(16q)-q/(16q)\ge7/8.
\]
The readout $2\sum_af_a\chi_a=\sum_af_at_a+\sum_af_a$
has margin at least $7/4$; implement its constant by the carried
$1/2$ rail. This is an ordinary-tanh network with two hidden
layers, polynomial width, polynomial weights, and polynomial product.

Attach it to the augmented prefix. The resulting network has
depth $j+3=O(\log(P+2))$ and polynomial width and product.
Some interface weights can be large if $S_j$ is tiny. Apply
Lemma~\ref{lem:physical} to this now-short network. It preserves
ordinary tanh and the Boolean function, makes all physical weights
polynomial in $P$, and rounds them with
$O(D'\log B)=O(\log^2(P+2))$ bits, since both the new depth
$D'$ and $\log B$ are $O(\log(P+2))$.
This proves all parts of the normal form.
\end{proof}


\begin{thebibliography}{9}
\bibitem{datta}
S.~Datta, M.~Mahajan, B.~V.~Raghavendra Rao, M.~Thomas, and H.~Vollmer.
\emph{Counting Classes and the Fine Structure between NC1 and L}.
2010. Primary full version:
\url{https://www.cse.iitm.ac.in/~bvrr/conf/mfcs10.pdf}.
In particular, Proposition~2.4 and Theorems~4.1--4.2.

\bibitem{caussinus}
H.~Caussinus, P.~McKenzie, D.~Th\'erien, and H.~Vollmer.
Nondeterministic NC1 computation.
\emph{Journal of Computer and System Sciences} 57(2):200--212, 1998.
\url{https://doi.org/10.1006/jcss.1998.1588}.

\bibitem{merrill}
W.~Merrill, H.~Jiang, Y.~Li, A.~W.~Lin, and A.~Sabharwal.
Why Are Linear RNNs More Parallelizable?
\emph{Proceedings of the 43rd International Conference on Machine Learning},
PMLR 306:88230--88260, 2026.
\url{https://proceedings.mlr.press/v306/merrill26a.html}.

\bibitem{unrestricted}
ChatGPT-6 Astra, K.~H\"anni, and D.~Vaintrob.
\emph{Representation constraints and polynomial-depth learning
in Gaussian Bayesian networks}, research draft, 19 October 2026.
The posterior ceiling used here is the theorem entitled
``A finite-scale posterior energy ceiling'' in
\path{works-in-progress/unrestricted-depth/source/unrestricted-primary-posterior-ceiling.tex}.
\url{https://github.com/kaarelh/AI-works/tree/main/neural-network-generalization/works-in-progress/unrestricted-depth}.

\bibitem{parallel}
ChatGPT-6 Astra, K.~H\"anni, and D.~Vaintrob.
\emph{Frobenius product, parallel circuits, and sample-complexity barriers
for deep Gaussian Bayesian networks}, 2026.
\url{https://github.com/kaarelh/AI-works/tree/main/neural-network-generalization/works-in-progress/boolean-depth-barriers}.
The two-ball geometry and two-majority statistical reduction are developed
there; the versions needed for the present theorem are proved here.
\end{thebibliography}
\end{document}
